\documentclass[a4paper,12pt]{article}

\usepackage[utf8]{inputenc} % encodé en utf-8
\usepackage[T1]{fontenc} % compatible avec les accents

\usepackage[french]{babel} % rédigé en français
\usepackage[hyphens]{url} % formatte les liens en autorisant la césure au niveau des traits d'union
\usepackage[urlcolor=black,colorlinks=true,linkcolor=black,citecolor=black]{hyperref} % liens cliquables mais non colorés

\usepackage[top=25mm, bottom=25mm, left=35mm, right=25mm]{geometry} % définit les marges
\usepackage{array} % permet l'alignement vertical dans les tableaux
\usepackage{setspace} % permet de régler l'interligne
\onehalfspacing % définit l'interligne 1,5 par défaut

\usepackage{graphicx} % gestion des images
\graphicspath{ {../images/}}
\usepackage{array} % gestion des tableaux
\usepackage{csquotes} % gestion des guillemets

\usepackage{fourier} % utilise une autre police que celle par défaut (Computer Modern) fourier - venturis - libertine - tgtermes - kpfonts - lmodern - txfonts - newtxtext,newtxmath - mathptmx - stix 

\usepackage[backend = biber, style = authoryear, useprefix = true, maxbibnames = 99]{biblatex}


%\addbibresource{}

\DefineBibliographyExtras{french}{%
\protected\def\mkbibnamefamily#1{%
\textsc{\textnohyphenation{#1}}}
}%

\DefineBibliographyStrings{french}{%
    page  = {\lowercase{p.}}, % for single page number
    pages = {\lowercase{pp.}}, % for multiple page numbers
}

\DefineBibliographyExtras{french}{\renewcommand*{\bibrangedash}{\textendash}} % Pour rendre les "--" correctement

% \renewcommand*{\mkibid}{\emph} %ibid en italique  if authoryear-ibid

\usepackage{float}

\setcounter{secnumdepth}{3} %Ne compte pas plus loin que 1.1.1
\setcounter{tocdepth}{3}

\usepackage{enumitem} %% utilisé pour la présentation de fausse table des matières dans la table des matières commentée

\newlist{legal}{enumerate}{10}
\setlist[legal]{label*=\arabic*.}

\usepackage{ragged2e} % Permet de placer plus facilement les tableaux

\newcolumntype{R}[1]{>{\RaggedRight}p{#1}}

\title{Projet STIC-B505 : Budgetsquirrel, une application de gestion de budget} 
\author{Milena Albu \and Marie Huysman} % vos prénom et nom
\date{} % pas de date

%\usepackage{fancyhdr}
%\pagestyle{fancy} % Turn on the style
%\fancyhf{}
%\fancyhead[L]{\footnotesize \nouppercase \leftmark}
%\fancyhead[R]{\footnotesize \nouppercase \rightmark}
%\fancyfoot[C]{\thepage}
%
%\renewcommand{\chaptermark}[1]{%
%\markboth{#1}{}}
%
%    \renewcommand{\sectionmark}[1]{\markright{\thesection \ #1}{}}

\usepackage{longtable}


\begin{document} % début du corps du texte
\sloppy

\begin{titlepage} % style spécial pour la page de garde
\singlespacing % remet l'interligne normal pour cette page

%%alternative à la bannière image
%\begin{center} 
%\huge{\textsc{Université Libre de Bruxelles}}\\
%\vspace{.5cm}
%\LARGE{Faculté de Lettres, Traduction et Communication}
%\end{center}
%% Bannière image
\begin{figure}[h]
\noindent
\makebox[\textwidth]{\includegraphics[scale=0.6]{banner.png}}
\end{figure}

\vfill % laisse LaTeX gérer l'espacement vertical

%VOIR SI LE TITRE CONVIENT

\begin{center}
\Huge{Budgetsquirrel}\\ % ajoutez votre titre ici
\LARGE{Application de gestion de budget - PHP et SQL}
\vspace{.5cm}
\end{center}

\vfill

\begin{tabular}{b{5.5cm}b{7.5cm}} % taille et alignement des cellules du tableau
 \\ Milena \textsc{Albu} \\ Marie \textsc{Huysman} \\ & Rapport de projet écrit dans le cadre du cours STIC-B505 : Conception et gestion de banques de données \\
\end{tabular}

\vfill

\begin{center}
Année académique 2020--2021 % vérifiez que l'année est correcte !
\end{center}

\end{titlepage}

% Nécessaire d'avoir une TOC pour un travail aussi court?
%\tableofcontents
\thispagestyle{plain}

\section{Introduction}

Si besoin de citer : \parencite[2--4]{SOURCE}
Notes de bas de page \footnote{informations suppl.}

\section{Schéma conceptuel}

\section{Schéma logique relationnel}

\section{Description du code de création de la base de données}

\section{Description de l'application Web}

\section{Conclusion}

\newpage

%voir si bibliographie nécessaire
%\chapter*{Bibliographie} 
%\addcontentsline{toc}{chapter}{Bibliographie}
%\markboth{Bibliographie}{}
%
%\printbibliography[heading=none]
%
%\fussy

%Voir si annexes nécessaire
%\chapter*{Annexes}
%\addcontentsline{toc}{chapter}{Annexes}
%\markboth{Annexes}{}

\end{document} % fin du corps du texte